\documentclass[12pt,a4paper]{article}
\usepackage[T1]{fontenc}
\usepackage[utf8]{inputenc}
\usepackage{lmodern,amsmath,amssymb,amsthm,mathtools}
\usepackage[a4paper,margin=25mm]{geometry}
\usepackage{microtype,booktabs,array,enumitem,xcolor}
\usepackage[colorlinks=true,linkcolor=blue!50!black,urlcolor=blue!50!black,citecolor=blue!50!black]{hyperref}
\setlength{\emergencystretch}{3em}
\setlength{\parskip}{.3em}
\setlist{nosep}
\newtheorem{theorem}{Theorem}
\newtheorem{lemma}[theorem]{Lemma}
\newtheorem{proposition}[theorem]{Proposition}
\newcommand{\HV}{\operatorname{HV}}
\newcommand{\Mag}{\operatorname{Mag}}
\newcommand{\wtO}{\widetilde O}
\newcommand{\1}{\mathbf1}
\newcommand{\SumOut}{\operatorname{SumOut}}
\newcommand{\Push}{\operatorname{Push}}
\hypersetup{pdftitle={Numerical Hypervolume and Magnitude in Dimensions Two to Ten},
pdfsubject={Finite weighted chains, low-dimensional sweeps, and pairwise Chan recursion}}
\title{Numerical Hypervolume and Magnitude\\in Dimensions Two to Ten\\[.5em]
\large Sweeps in 2D and 3D; pairwise Chan recursion in 4D--10D}
\author{Research working report}
\date{4 October 2026}
\begin{document}
\maketitle
\begin{abstract}
We extend the numerical four-dimensional construction to every integer
dimension from two to ten, including odd dimensions. The input is a finite
set of nonnegative corners defining boxes anchored at the origin. The same
coordinate-cell representation computes ordinary hypervolume and the
dominated-set $\ell_1$ magnitude by changing the one-dimensional masses.
No symbolic integration or numerical quadrature is used.

Two and three dimensions use direct numerical sweeps with
$O(n\log n)$ arithmetic complexity. In dimensions four through ten, a
dimension-independent prefix-sum elimination kernel, paired coordinate
order, weighted geometric cuts, and block-mass compression give
$\wtO_d(n^{d/3})$ arithmetic complexity. Odd dimensions finish each round
with one singleton operation. Compression temporarily uses $2d$ variables,
so the ten-dimensional implementation includes twenty-variable compression.
An explicit local block schedule accounts for representation growth.

The accompanying standalone Python implementation supports all nine
dimensions. Exact inclusion--exclusion tests, skyline-update tests, and
compression tests accompany it. The exponent is Chan's; this is a numerical
formulation and implementation, with potentially large dimension-dependent
constants, not a claim of improved exponents or competitive practical speed.
\end{abstract}
\newpage
\tableofcontents
\newpage

\section{What is now available}
The implementation \texttt{numerical\_chan.py} accepts dimensions
$d=2,3,\ldots,10$. The earlier four-dimensional file remains a separate
reference implementation. A user-facing call is
\begin{quote}\ttfamily\small
from numerical\_chan import hypervolume, magnitude, NumericalChan\\
hv = hypervolume(points)\hfill\# infer dimension\\
mag = magnitude(points)\\
solver = NumericalChan(dimension=7)\\
hv = solver.compute(points)
\end{quote}
Here \texttt{points} contains nonnegative box corners; it is not a list
of arbitrary lower and upper box endpoints. A reference-point conversion
for conventional minimization data is given in Section~\ref{sec:usage}.

\begin{theorem}[Fixed-dimensional numerical bounds]\label{thm:main}
For every fixed $d\in\{2,\ldots,10\}$, the algorithm computes the
hypervolume and dominated-set $\ell_1$ magnitude of $n$ anchored boxes.
Its arithmetic complexity is
\[
\begin{cases}
O(n\log(n+2)),&d=2,3,\\
O_d(n^{d/3}\log^{c_d}(n+2)),&4\le d\le10,
\end{cases}
\]
where $c_d$ is a constant depending only on $d$. All evaluations use finite
numerical arrays, comparisons, integer searches, and arithmetic. No
Cartesian product grid of coordinate cells is constructed by the solver.
\end{theorem}
Arithmetic has unit cost in this theorem. Exact rational arithmetic can
have additional bit cost. Floating-point use can suffer cancellation;
the theorem is not a numerical-error guarantee. The notation $\wtO_d$
suppresses polylogarithms in $n$ and constants depending on the fixed
dimension. It does not suppress polynomial factors in $n$.

\begin{center}
\begin{tabular}{@{}clcc@{}}
\toprule
$d$&Numerical method&Bound&Variables in compression\\\midrule
2&Descending sweep&$O(n\log n)$&---\\
3&Sweep with a skyline tree&$O(n\log n)$&---\\
4&$2+2$ recursion&$\wtO(n^{4/3})$&8\\
5&$2+2+1$ recursion&$\wtO(n^{5/3})$&10\\
6&$2+2+2$ recursion&$\wtO(n^2)$&12\\
7&$2+2+2+1$ recursion&$\wtO(n^{7/3})$&14\\
8&$2+2+2+2$ recursion&$\wtO(n^{8/3})$&16\\
9&$2+2+2+2+1$ recursion&$\wtO(n^3)$&18\\
10&$2+2+2+2+2$ recursion&$\wtO(n^{10/3})$&20\\
\bottomrule
\end{tabular}
\end{center}
The pair notation specifies elimination and cut order. It does not assert
that the geometric problem factors into independent two-dimensional
problems. Relations between different pairs are retained.

\section{The common numerical representation}
For $P\subset\mathbb R_{\ge0}^d$, let
\[
D(P)=\bigcup_{p\in P}\prod_{i=1}^d[0,p_i].
\]
We compute either $\HV_d(P)=\lambda_d(D(P))$ or
\[
\Mag(D(P))=\left(\bigotimes_{i=1}^d
           (\delta_0+\tfrac12\lambda)\right)(D(P)),
\]
using Emmerich's dominated-set magnitude formula~\cite{emmerich}.
The measure includes mass on coordinate faces and at the origin.

On each coordinate axis, sort the distinct positive corner coordinates:
$0=c_{i,0}<c_{i,1}<\cdots<c_{i,m_i}$. Use cells $\{0\}$ and
$(c_{i,k-1},c_{i,k}]$. Their numerical masses are
\begin{equation}
\begin{array}{c|cc}
&a_i(0)&a_i(k),\ k\ge1\\\hline
\text{hypervolume}&0&c_{i,k}-c_{i,k-1}\\
\text{magnitude}&1&(c_{i,k}-c_{i,k-1})/2
\end{array}                                      \label{eq:masses}
\end{equation}
Replace each corner by its rank vector $r(p)$. On cell indices $k$,
box membership is exactly $k_i\le r_i(p)$ for all $i$.
Consequently, summing the product masses of covered index cells is exact.
Only the $d$ axis arrays are constructed, with total length at most
$d(n+1)$; sorting and rank lookup cost $O_d(n\log n)$.

The anchor cell is retained even when it has zero hypervolume mass.
The empty input has value zero for both measures, whereas a single
origin point has hypervolume zero and magnitude one. Duplicate corners
and tied coordinates require no general-position assumption.

\section{Two dimensions: one descending sweep}
Let $A_2(t)=\sum_{j=0}^t a_2(j)$, and set $A_2(-1)=0$.
Scan the first-axis index $x$ in decreasing order. Before charging cell
$x$, insert all boxes whose first coordinate rank is $x$.
Maintain
\[
H=\max\{A_2(r_2(p)):r_1(p)\ge x\},
\]
with empty maximum zero. Add $a_1(x)H$ to the answer.
All intervals in the second coordinate share their anchor, so their union
mass is the largest prefix mass. Each point is inserted once.

The scan itself is linear after sorting. Including preprocessing gives
$O(n\log n)$ time and $O(n)$ space. This computes either measure simply
by using the appropriate masses in \eqref{eq:masses}.

\section{Three dimensions: a numerical skyline tree}\label{sec:3d}
Again scan first-axis indices in decreasing order, inserting every box
before charging its first-axis cell. The active cross-section is a union
of anchored rectangles in coordinates two and three.
Let $A_3(t)=\sum_{j=0}^t a_3(j)$, with $A_3(-1)=0$, and define
\[
h(y)=\max\{A_3(r_3(p)): p\text{ is active and }r_2(p)\ge y\},
\]
where the empty maximum is zero. The cross-section mass is
\begin{equation}
                         V=\sum_y a_2(y)h(y).       \label{eq:skyline}
\end{equation}
The array $h$ is nonincreasing. Inserting a rectangle with second rank $u$
and third-coordinate prefix mass $v$ performs the update
\[
                         h(y)\gets\max(h(y),v)\quad(y\le u).
\]

\begin{lemma}[One interval changes]
Let $\ell$ be the first index with $h(\ell)<v$, or the array length if
there is none. If $\ell>u$, nothing changes. Otherwise the update assigns
the constant value $v$ on exactly $[\ell,u]$.
\end{lemma}
\begin{proof}
Nonincreasing order implies $h(y)\ge v$ before $\ell$, and $h(y)<v$
at every subsequent index. The insertion affects only the prefix ending
at $u$. Its changed set is therefore the stated interval. The updated
array remains nonincreasing at both ends of that interval.
\end{proof}

A lazy segment tree stores, for each index interval $I$, its minimum
height and its weighted total $\sum_{y\in I}a_2(y)h(y)$. It also stores
a pending constant assignment when one exists. The mass
$\sum_{y\in I}a_2(y)$ is obtained from a fixed prefix array.
\begin{enumerate}
\item Find the first height below $v$ by descending into the first tree
child whose minimum is below $v$.
\item If it lies in $[0,u]$, assign $v$ on $[\ell,u]$ with lazy propagation.
An assigned node's weighted total becomes $v\sum_{y\in I}a_2(y)$.
\item Read $V$ at the tree root, and add $a_1(x)V$ after all insertions at $x$.
\end{enumerate}
Both tree operations cost $O(\log n)$, and there are $n$ insertions and
$O(n)$ first-axis cells. This proves the three-dimensional part of
Theorem~\ref{thm:main}, with $O(n)$ space.
Zero-weight cells cause no exception: comparisons concern heights,
and a weighted sum can include zero masses.

\section{The common elimination kernel for dimensions four to ten}
This section gives the same finite-array kernel for any fixed number
$p$ of variables. During compression, $p$ can be as large as $2d$.
Each variable $k_i$ ranges over a finite chain
$E_i=\{0,\ldots,s_i-1\}$. A term is
\begin{equation}
T(k)=c\prod_i a_i(k_i)\prod_e\1[k_{j_e}\mathrel{\bowtie_e}f_e(k_{i_e})],
\qquad i_e<j_e,\quad\bowtie_e\in\{\le,\ge\}. \label{eq:term}
\end{equation}
The unary arrays and scalar coefficient can be signed. Cutoff arrays are
integer-valued and monotone. For each unordered pair retain at most four
slots: upper or lower bound, each with nondecreasing or nonincreasing
cutoff. Within one slot, combine upper bounds by minimum and lower bounds
by maximum. Opposite monotonicities remain separate.

Thus a term uses $O_p(M)$ entries, where $M=\sum_i s_i$.
Sentinels $-1$ and $s_j$ encode empty upper and lower rays; no infinite
numerical values or decorated open endpoints are needed.

\subsection{Comparisons remain numerical}
An active bound is either a constant or a monotone array $A(k_i)$.
To impose $A(k_i)\le B(k_j)$, use the following rules.
\begin{itemize}
\item If the two variables coincide, multiply their unary array by the
Boolean comparison array. It may oscillate; unary weights need not be monotone.
\item If one bound is constant, the comparison is also a unary Boolean array.
\item For distinct variables, invert the monotone array on the larger-index
variable by binary search. This produces one monotone cutoff in a pair slot.
\end{itemize}
Strict comparisons use the corresponding strict search boundary. The
total work per comparison is $O_p(M\log(M+2))$.

\subsection{Eliminating one variable}
For an index $q\in\{0,\ldots,s-1\}$ to eliminate, invert every incident
pair predicate into a lower or upper bound on $q$. Include constant bounds
$0$ and $s-1$. At most $2p-1$ lower and $2p-1$ upper candidates occur,
because each other variable supplies at most two of each kind.

For each pair of candidate labels, impose the comparisons saying that
$L$ is the first maximizing lower bound and $U$ the first minimizing
upper bound, and impose $L\le U$. Earlier labels are defeated strictly;
later labels are defeated weakly. These regions are disjoint and exhaust
the feasible assignments of the surviving variables.

Build the numerical prefix $P(t)=\sum_{q<t}a(q)$. On a label region,
the eliminated sum is exactly
\begin{equation}
                            P(U+1)-P(L).            \label{eq:prefix}
\end{equation}
Each summand is a scalar or a unary array in the variable supplying the
bound. Multiply it into the surviving term with the appropriate sign.
When both bounds use the same variable, combine their difference directly.
The prefix need not be monotone: it becomes a unary weight, never an
inverse cutoff.

\begin{lemma}[Closure]
One application of $\SumOut$ emits at most
$B_p=2(2p-1)^2$ terms in the same four-slot class, in
$O_p(M\log(M+2))$ arithmetic operations per input term.
\end{lemma}
\begin{proof}
There are at most $(2p-1)^2$ winning label pairs. Each emits at most two
terms from \eqref{eq:prefix}. All label comparisons are of the forms just
described; merging same-slot constraints preserves monotonicity. No new
variable or new coordinate domain is introduced. A dimension-bounded
number of array scans and binary searches gives the work bound.
\end{proof}

\subsection{Pairs and the odd-dimensional final step}
For an ordered axis list $(1,\ldots,d)$, apply $\SumOut$ to $(1,2)$,
then $(3,4)$, and so on. If $d$ is odd, eliminate $d$ alone at the end.
This is exactly the same one-variable operation, not a new odd-dimensional
case or a padding dimension. Cross-pair constraints remain in the state.
No Cartesian table on a coordinate pair is constructed.

\section{Compression by block masses in $2d$ variables}\label{sec:push}
Suppose a geometric cell has $m$ surviving hard boxes. On each current
axis, place a block boundary immediately after each surviving box cutoff
strictly below the cell's upper endpoint. The blocks are contiguous integer
intervals $I_i(t_i)=[L_i(t_i),R_i(t_i)]$, with at most $m+1$ per axis.
Define the desired coarse mass by
\begin{equation}
        (\Push F)(t)=\sum_{k_i\in I_i(t_i),\ 1\le i\le d}F(k).
                                                        \label{eq:push}
\end{equation}
This equation specifies the answer; it does not instruct the implementation
to enumerate the block products.

For each old term, add new variables $t_1,\ldots,t_d$ with unit unary
weights and the $2d$ membership comparisons
$L_i(t_i)\le k_i\le R_i(t_i)$. There are now $2d$ variables.
Eliminate only the old variables, in consecutive pairs with an odd final
singleton when necessary. Rename the $d$ surviving block variables.

\begin{lemma}[Mass compression and reuse]\label{lem:push}
If the old state has $S$ terms and total axis length $M$, compression
emits at most $C_d S$ terms, where
\[
                           C_d=\prod_{p=d+1}^{2d}B_p.
\]
It costs $O_d(S(M+m)\log(M+m+2))$ operations. The output uses
$O_d(C_d S(m+1))$ array entries. For any future mask $g$ constant on
the current block products,
\[
                      \sum_k F(k)g(k)=\sum_t(\Push F)(t)g(t).
\]
\end{lemma}
\begin{proof}
Apply the closure lemma $d$ times. The combined old and new axis length is
$M+O_d(m+1)$ and only decreases during elimination. Every remaining array
is indexed by one coarse axis. For the final identity, partition the
finite sum by block products and move the constant value of $g$ outside
each inner sum. Every surviving box cutoff was used to create the blocks;
future box masks and cuts therefore have the required constancy.
\end{proof}
No density is reconstructed and no division by cell measure is performed.
An anchor may be merged into a neighbouring block when future predicates
do not distinguish it; its mass is then retained in the block sum.
For $d=10$ this operation genuinely lifts to twenty variables and removes
ten. The implementation does not substitute a four-dimensional compressor.

\section{The geometric recursion in any fixed $d\ge4$}
The recursion stores the current index cell $\Gamma=\prod_i[\ell_i,u_i]$,
a numerical state $F$, and the list $B$ of unabsorbed boxes. Its invariant is
\begin{equation}
U(F,B,\Gamma)=\sum_{k\in\Gamma}F(k)
                   \prod_{p\in B}(1-\1[k\le p]).   \label{eq:invariant}
\end{equation}
Initially $F=\prod_i a_i$. Return total bounding-box mass minus $U$.
After compression, $F$ stores coarse masses; the same invariant holds.

\subsection{Absorb one- and two-sided boxes}
Discard a box if some $p_i<\ell_i$. Otherwise call axis $i$ active if
$p_i<u_i$. A box with no active axis covers the cell, so $U=0$.
A box with one active axis $i$ is absorbed by zeroing the corresponding
unary entries $k_i\le p_i$ in every term.

For a group of boxes with active axes $i<j$, build the suffix maximum
\[
f(k_i)=\max\{p_j:p_i\ge k_i\},\qquad\max\varnothing=-1.
\]
Their common uncovered predicate is $k_j\ge f(k_i)+1$, a nonincreasing
lower cutoff. Absorb it into the state and discard the group's boxes.
The suffix array is built for the group, not separately for every box.
Absorption never increases the number of terms. Every retained hard box
has at least three active axes.

\subsection{Weighted cuts and one full round}
Cycle through axes $1,\ldots,d$. Relative to the current axis, their
cyclic positions are $0,\ldots,d-1$. Each triple $J$ of active axes of
each hard box contributes weight
\[
                            \omega(J)=2^{\sum_{j\in J}\mathrm{pos}(j)/d}.
\]
Let $\Phi$ be their sum. For fixed $d$, $\Phi$ is within constant factors
of the number of hard boxes. On the current axis, take a weighted median
of the cutoff coordinates of triples containing that axis.
Split after that index $c$ into $[\ell,c]$ and $[c+1,u]$.
These cells are disjoint, including at atomic anchors.

\begin{lemma}[Potential contraction]
After the cut and cyclic advance, each child has potential at most
$2^{-3/d}\Phi$. If no triple uses the axis, skip the cut and advance;
the same bound holds.
\end{lemma}
\begin{proof}
At most half of the weight of triples using the cut axis survives in
either child. A surviving such triple increases its position sum by
$d-3$, so its weight changes by $2^{(d-3)/d}$; half of this factor is
$2^{-3/d}$. Every other triple reduces its position sum by three and
changes by exactly $2^{-3/d}$. Restriction introduces no new active triple.
If there is no candidate, only the latter case is present.
\end{proof}
It follows that a full $d$-axis round has at most $2^d$ descendants,
each with potential at most $\Phi/8$. This explains the exponent $d/3$,
but a total-work proof must also charge the numerical state.

\subsection{Integer weights and the block schedule}
With $n_0$ original points, set $A=(n_0+2)^2$. A triple's position sum
$s$ is an integer from $3$ through $3d-6$. Use
\[
                     W_s=\left\lceil(A^d2^s)^{1/d}\right\rceil
\]
as its integer median weight. Binary search computes this root using
integer powers; no floating-point root is used. The approximation satisfies
\[
        2^{s/d}\le W_s/A\le(1+1/A)2^{s/d}.
\]
Thus the contraction gains at most a factor $1+1/A$ per level. Over
$O_d(\log n_0)$ levels the total distortion is bounded by a constant.
The combinatorial weights use $O_d(\log n_0)$ bits.

At the start of each block of geometric levels, set
\begin{equation}
 Q=\left\lceil\frac{\sum W_s}{A}\right\rceil,\qquad
 q=\max\left\{1,\left\lfloor\frac{\log_2 Q}{4d}\right\rfloor\right\}.
                                                        \label{eq:schedule}
\end{equation}
Run $q$ full rounds, or $dq$ levels. At each surviving block leaf absorb
newly easy boxes, compress using the remaining hard cutoffs, and start
a new block using its own $Q$. Skips consume one level just as cuts do.
For odd $d$, a round consists of pairs and the last singleton axis.

\subsection{Explicit recursive procedure}
\begin{quote}\ttfamily\small
Solve(state, boxes, cell, axis, levels-left):\\
\hspace*{1em}discard disjoint boxes; absorb all easy boxes\\
\hspace*{1em}if a box covers cell or state is zero: return 0\\
\hspace*{1em}if at most b hard boxes remain: return Base(state, boxes, cell)\\
\hspace*{1em}if levels-left is zero:\\
\hspace*{2em}state = Push(state, remaining hard-cutoff blocks)\\
\hspace*{2em}relabel boxes and cell; initialize dq levels by (\ref{eq:schedule})\\
\hspace*{1em}choose a weighted median on axis\\
\hspace*{1em}if none: recurse on same cell with next axis and one fewer level\\
\hspace*{1em}else: recurse on both disjoint children and add their answers
\end{quote}
The first block is also initialized by \eqref{eq:schedule}. For a fixed
constant $b$ (default two), the base case expands the product in
\eqref{eq:invariant} into $2^b$ signed rectangle sums. Evaluate each
rectangle by paired $\SumOut$ calls, with an odd final singleton.
This inclusion--exclusion is constant-size; it is not applied to the full
input list.

\section{Proof of the full arithmetic bound}
Discretization, absorption, disjoint splitting, exact one-variable sums,
and Lemma~\ref{lem:push} prove correctness inductively. The base case is
the finite inclusion--exclusion identity, so the returned value is exact.

For the work bound, let $N=\Phi$ and $S$ be the state term count at a
block entry. Apart from the root, the axis length is $O_d(N)$ because
the preceding compression retained only surviving hard cutoffs. The root
has $O_d(n)$ entries and can be charged at root size even if many boxes
are absorbed immediately.

Put $\beta_d=1-3/(4d)$. For sufficiently large $N$,
\eqref{eq:schedule} gives $q=(\log_2N)/(4d)+O_d(1)$.
There are $O_d(2^{dq})=O_d(N^{1/4})$ nodes and leaves in the block.
Each is charged $\wtO_d(SN)$ work, including compression or an early
numerical base case. Each surviving child block has potential
$O_d(N^{\beta_d})$ and at most $C_dS$ terms. Thus
\begin{equation}
T_d(N,S)\le c_1N^{1/4}T_d(c_2N^{\beta_d},C_dS)
                         +\wtO_d(SN^{5/4}).        \label{eq:recurrence}
\end{equation}
All constants depend only on $d$. Once $N$ is bounded by a fixed constant,
only a bounded number of further levels is possible by geometric
contraction, so the base cost is $O_d(S)$.

At macro-depth $g$, the largest size is at most $c_3N^{\beta_d^g}$.
There are $O_d(\log\log(N+2))$ macro-generations, giving at most
$C_d^gS=S\log^{O_d(1)}(N+2)$ terms along a path.
No assumption that empirical merging keeps the term count small is needed.
The number of macro-nodes at depth $g$ is at most
\[
c_4^gN^{\frac14\sum_{h<g}\beta_d^h}
       =c_4^gN^{\frac d3(1-\beta_d^g)}.
\]
Multiplying by the per-node charge gives, apart from polylogarithms,
\[
 S N^{\frac d3(1-\beta_d^g)+\frac54\beta_d^g}
 =S N^{\frac d3-(\frac d3-\frac54)\beta_d^g}
 \le S N^{d/3},
\]
because $d/3>5/4$ for $d\ge4$. Constants raised to $g$ and the number
of levels contribute only polylogarithmic factors. Summation proves
$T_d(N,S)=\wtO_d(SN^{d/3})$. At the root $S=1$, $N=O_d(n)$; its
possibly longer arrays cost at most $\wtO_d(n^{5/4})$, also within the
bound. Together with the sweeps this proves Theorem~\ref{thm:main}.

This particular block recurrence would not establish $\wtO(n)$ for
$d=3$: its block-work exponent is $5/4$. That is why the implementation
selects the skyline sweep there. Merely substituting $d=3$ in the final
formula would be unjustified.

\section{Magnitude and hypervolume in both directions}
For any dimension covered here, finite inclusion--exclusion and
coordinatewise minima give
\begin{equation}
            \Mag(D(P))=\HV_d(\{\1+p/2:p\in P\}).   \label{eq:affine}
\end{equation}
Indeed, the mass of one anchored box is $\prod_i(1+p_i/2)$, and the
affine map commutes with intersections described by coordinatewise minima.
The identity also holds for empty input.

Conversely, if every coordinate of every HV corner $q$ is at least one,
\[
                    \HV_d(Q)=\Mag(D(\{2(q-\1):q\in Q\})).
\]
For arbitrary strictly positive corners choose $s>0$ with $sq_i\ge1$;
then
\[
                  \HV_d(Q)=s^{-d}\Mag(D(\{2(sq-\1):q\in Q\})).
\]
Zero-coordinate boxes can first be discarded for this inverse HV
transformation because their $d$-dimensional volume is zero. They must
not be discarded for direct magnitude computation. If no boxes remain,
return zero. These transformations preserve the polynomial exponent.

\section{Examples and use of the implementation}\label{sec:usage}
The same two-box example works in every supported dimension:
\[
p=(2,1,1,\ldots,1),\qquad q=(1,2,1,\ldots,1).
\]
Their intersection has every extent one. Hence
\[
               \HV_d(\{p,q\})=3,\qquad
               \Mag(D(\{p,q\}))=\frac{15}{4}(\tfrac32)^{d-2}.
\]
An example call in seven dimensions is
\begin{quote}\ttfamily\small
points = [(2, 1, 1, 1, 1, 1, 1), (1, 2, 1, 1, 1, 1, 1)]\\
assert hypervolume(points) == 3\\
assert magnitude(points) == Fraction(3645, 128)
\end{quote}
Here \texttt{Fraction} is imported from the standard-library
\texttt{fractions} module. Integer and \texttt{Fraction} coordinates retain
exact rational arithmetic. No third-party Python package is needed.

For minimization points $a$ and a reference point $r$ with $a_i\le r_i$,
convert to extents $p_i=r_i-a_i$. Then
$\lambda_d(\bigcup_a[a,r])=\HV_d(\{r-a\})$ by reflection and translation.
Input points outside this reference region must be handled explicitly;
the core interface rejects negative extents. The magnitude interface
applies to the resulting origin-anchored dominated set and its product
measure, as defined above.

\texttt{NumericalChan(dimension=d)} fixes the expected dimension and
checks every point. Without that argument, a nonempty input supplies its
dimension; an empty input returns zero. All dimensions must lie between
two and ten. Nonfinite or negative coordinates, mixed coordinate counts,
and out-of-range dimensions are rejected.

The optional \texttt{block\_levels} parameter is a test override. An
arbitrary fixed override is not the schedule covered by the asymptotic
theorem. The low-dimensional sweeps do not use it. Production calls should
leave it unset. The implementation exposes counters through
\texttt{solver.stats} for diagnostics.

\section{The Numba backend and measured performance}\label{sec:numba}
The file \texttt{numerical\_chan\_numba.py} provides a separate
\texttt{float64} backend. In dimensions two and three, the numerical sweep
and skyline-tree operations run in Numba's nopython mode. In dimensions
four through ten, Numba compiles array scans, monotonicity checks, prefix
sums, cutoff searches, normalization, and suffix masks. Geometric recursion,
winning-bound enumeration, and structural term management remain Python.
This is a hybrid implementation. No fastmath or parallel execution is used.

The exact standard-library backend remains available independently.
The Numba version converts its coordinates and masses to floating point;
it does not preserve exact rational arithmetic. A first call may include
compilation or cache loading. Reported steady-state timings exclude that
first call, and separate fresh-cache probes record first-use cost.

The reproducible benchmark compares the numerical Python and Numba
backends with the repository's unmodified Python implementations of Chan
Section~4.2 ($d/3$) and Section~2 ($d/2$). The baseline commit is
\texttt{08539aa99e94}; neither baseline uses Numba. All timed inputs are
floating point. Dominance prefiltering is disabled, and input preprocessing
and coordinate transforms are included. For magnitude, the references use
the affine HV transformation.

Each entry below is the median of three calls after one untimed first call.
Workers run sequentially. A 35-second limit covers the entire worker,
including startup, its first call and all repeats; ``timeout'' is not a
per-call runtime. Raw repeats and first-call times are in the accompanying
JSON/CSV. No ratio is assigned to a timeout.

\begin{center}\small
\begin{tabular}{@{}lrrrr@{}}
\toprule
Case&Numerical Python&Numba&Chan 4.2&Chan 2\\
\midrule
2D, $n=256$&0.0014577&0.001104&---&0.013813\\
3D, $n=128$&0.0017669&0.0006409&0.047406&0.015749\\
4D, $n=24$&1.4637&2.3398&0.027249&0.0054424\\
4D, $n=48$&timeout&timeout&0.10325&0.018492\\
5D, $n=16$&1.1589&2.2592&0.049136&0.0066998\\
6D, $n=12$&1.3261&2.8731&0.060131&0.0075414\\
7D, $n=8$&0.75676&1.7606&0.055029&0.0070965\\
8D, $n=10$&7.6582&timeout&0.19733&0.03268\\
9D, $n=6$&timeout&timeout&0.03161&0.0032998\\
10D, $n=8$&timeout&timeout&0.20869&0.023293\\
\bottomrule
\end{tabular}
\end{center}
All times are seconds. The ND overview uses spherical, mutually
nondominated fronts; the full data additionally includes tied-grid inputs
and magnitude evaluations. These small cases do not measure asymptotic exponents.

For completed matched ordinary-HV cases in dimensions 4--10, the median ratio Numba numerical / Python numerical was 1.93x (20 cases). The median ratio Numba numerical / the repository's Python Section~4.2 implementation was 21.25x (20 cases). Ratios above one mean the Numba version was slower. These are medians of per-case ratios for this small benchmark set, not universal performance predictions. Array-level JIT compilation does not remove the cost of enumerating and managing many terms. The results do not support a high-dimensional speedup claim for this implementation.

Across completed comparisons the maximum scaled disagreement was $3.31\times10^{-16}$, where scaling divides by $\max(1,|\text{reference value}|)$. This is an observed check, not a uniform error bound.

Fresh-cache first-use probes (module imports excluded) gave:
\begin{center}\small
\begin{tabular}{@{}crr@{}}\toprule
Dimension&First use (s)&Second use (s)\\\midrule
2&2.5352&9.72e-05\\
3&2.496&0.0001481\\
4&1.9391&0.0030841\\
\bottomrule
\end{tabular}
\end{center}
The first call includes both compilation and computation; it is not pure compiler time.

Environment: CPython 3.12.14, NumPy 2.5.3, Numba 0.68.0, Windows.


\section{Validation and practical scope}\label{sec:validation}
The verifier \texttt{verify\_numerical\_chan.py} uses independent exact
oracles. In each of the nine dimensions it checks empty input, the origin,
zero coordinates, duplicate and dominated corners, cyclic antichains,
random positive integer instances, and rational coordinates against
inclusion--exclusion. It also checks both directions of the affine
magnitude--HV transformation.

The skyline tree is compared after each insertion with a dense
one-dimensional array. Compression is checked at every lifted dimension
$2d=4,6,\ldots,20$: the verifier explicitly enumerates small fine grids
and compares their block sums with the emitted numerical state. It then
applies a new aligned cross-axis mask and compresses a second time.
The Cartesian enumeration belongs only to the independent verifier.

For dimensions four through ten, additional instances force repeated
geometric compression with a deliberately short test-only block length.
These extend a genuinely four-dimensional instance by common extents in
the additional coordinates, so the expected multiplicative factor is
known. These structured tests exercise high-dimensional execution; they
do not establish practical behaviour on large, fully coupled 10D inputs.
Small general instances are tested separately.

The exact integer-root inequalities are checked directly. Geometric
potential diagnostics use floating evaluation of the theoretical
algebraic weights with a tolerance; measure comparisons and compression
comparisons use exact arithmetic. Full counts and diagnostic states are
recorded in \texttt{numerical\_chan\_results.json}.
The exact suite passes 234 end-to-end cases, 248 first-compression cell
checks, 18 second compressions, and 14 forced repeated-recursion instances.
The separate Numba suite passes 144 end-to-end comparisons with an exact
inclusion--exclusion oracle, 2,044 compression-cell comparisons,
nine signed four-slot sums, six repeated-compression instances, and
320 direct checks of the compiled skyline update. Its maximum scaled
end-to-end discrepancy in this suite is $3.60\times10^{-14}$.
These are observed validation results, not a uniform error bound.

The constants $C_d$ can be very large and the hidden logarithmic exponent
depends on them. The prototype uses exact-key merging and normalization
of unary factors to reduce practical state size, but the correctness and
raw-emission proof do not require successful merging. If deterministic
comparison-based term handling is desired, optional merging can be omitted
or replaced by sorting. Even quadratic merging overhead in the
polylogarithmic term count changes only a hidden logarithmic factor.
These measurements do not establish an asymptotic speed advantage.

\section{Relation to Chan and the supplied repositories}
The geometric potential, absorption of locally one- and two-sided
orthants, and periodic compression follow Chan's orthant framework
in Sections 4.1--4.2 of~\cite{chan}. The simplification is its finite
numerical realization: explicit cell masses, four monotone slots per
coordinate pair, prefix-sum elimination, and block-mass pushforward.
The locally chosen blocks make representation growth part of the proof.

The supplied \texttt{FastHVChan} repository~\cite{fastrepo} and
\texttt{HV4DMagnitude} repository~\cite{magrepo}, including its
\texttt{HV468DMagnitude} reports, supplied the implementation context.
The earlier four-dimensional report described six- and eight-dimensional
extensions mathematically; the present implementation and tests extend
the actual numerical engine through ten dimensions, including five,
seven, and nine. Dedicated 2D and 3D sweeps complete the requested range.

The construction preserves Chan's high-dimensional exponent. It neither
proves a sub-$n^{4/3}$ four-dimensional bound nor requires the previously
unproved global batching lemma for that stronger target.

\begin{thebibliography}{9}
\raggedright
\bibitem{chan}
T. M. Chan. \emph{Klee's Measure Problem Made Easy}.
Proceedings of FOCS 2013, pp. 410--419. DOI: 10.1109/FOCS.2013.51.
Author's August 2013 version, Sections 4.1--4.2:
\url{https://tmc.web.engr.illinois.edu/easyklee8_13.pdf}.

\bibitem{emmerich}
M. T. M. Emmerich. \emph{The Magnitude of Dominated Sets: A Pareto Compliant
Indicator Grounded in Metric Geometry}. arXiv:2604.18147, 2026.
\url{https://arxiv.org/abs/2604.18147}.

\bibitem{fastrepo}
M. Emmerich and collaborators. \emph{FastHVChan}.
Repository snapshot \texttt{08539aa99e94}, inspected 4 October 2026.
\url{https://github.com/emmerichmtm/FastHVChan}.

\bibitem{magrepo}
M. Emmerich and collaborators. \emph{HV4DMagnitude}, including the
pair-elimination reports and the \texttt{HV468DMagnitude} code and audits.
Repository snapshot \texttt{0c668e765202}, inspected 4 October 2026.
\url{https://github.com/emmerichmtm/HV4DMagnitude}.
\end{thebibliography}
\end{document}
